% Transcription — fonds Alexandre Grothendieck, shelfmark 72, batch 1 (pages 1-20).
% Established from the facsimile archives/batches/72/batch-01.pdf (cut from a
% copy of 72.pdf fetched through the project's relay,
% https://grothendieck-relay.onrender.com/source/72.pdf, a mirror of
% https://grothendieck.umontpellier.fr/72.pdf), per the skill
% transcribe-grothendieck. The folder is 112 pages in six batches (1-20,
% 21-40, 41-60, 61-80, 81-100, 101-112), transcribed in parallel; this is the
% first. The raw scan's cover sheet was removed from the batch PDF: PDF page k
% is archivists' page k, and the pencilled numbers bottom-left (« 1 » ...
% « 20 ») agree.
%
% Pass: Opus 5.5 (claude-opus-5-5), 2026-09-27 — first pass, unchecked against the pages by a human.
%
% Pages skipped: none as such. Page 1 is a wrapper sheet inscribed in his
% hand, top right, « n-polyèdres réguliers et représentations de G_n dans
% O(n+1) ..... (dictionnaire) » (the sign in « n+1 » doubtful), with « (23 p) »
% in pencil to its left (whose hand is not settled). By the rule for covers it
% takes no \page marker; his words become the section heading, with a \note.
% Pages 9, 10, 12, 18 and 20 are written on their upper part only.
%
% Pages summarised: none. No page of the batch is typed or another's.
%
% His pagination: none visible on these leaves; no date in his hand.
%
% What the batch is. A « dictionnaire » between regular n-polytopes
% (« n-hyperpolyèdres combinatoires épinglés », with automorphism group
% 𝔊_n generated by τ_0, ..., τ_n, τ_i and τ_j commuting for i, j non
% consecutive — the Coxeter group of a string diagram) and their geometric
% realisations over a base scheme S, described in a string of equivalent
% categories:
%   2-4   the list C_pr, C'_pr, C_aff, C'_aff, C_vect, C'_vect, C^∞_pr,
%         C^∞_vect, C^∞_pr.orth, C^∞_vect.orth (projective, affine, vector,
%         « at infinity », orthogonal variants), which are self-dual, the
%         rigidifications, and the chain of equivalences on p. 4.
%   5-9   kernels of Aut(E, f, π, π') -> Aut(V) and -> Aut(C = P(V)): over
%         an integral ring / a field the only extra automorphism is the
%         « antipodisme » (in characteristic 2 or for ε = -1).
%   10    duality: C and C' in duality when ρ is invertible; exchanging
%         τ_i and τ_{n-i}.
%   11-12 twisting: realisations of a given Π, or of combinatorial polytopes
%         of type 𝔊_0 twisted by a torsor.
%   13-17 a fresh start: C_0 (projective realisations = triples (X, τ_*,
%         s_0)), C'_0 (vector bundle V with its canonical basis a_0..a_n,
%         a_i^∨ in terms of λ_i, β_i), C''_0 (C = P(V), equivalent if
%         n >= 2), and the invariant quadratic form f_V computed
%         explicitly on p. 17.
%   18-20 C'''_0 with the invariant hyperquadric Q_C; conditions b'), c');
%         case n even, G -> SO(f_V); example n = 2 with Y = Q_C a conic
%         bundle and the « dihedral » case excluded.
%
% Notation, fixed once. His gothic G (the group of the universal pinned
% polytope) is \mathfrak{G} (\mathfrak{G}_n where he writes the index); his
% plain G (p. 11, 13 ff.) is kept as G. « ps. réfl. » = pseudo-réflexion;
% « rg » = rang; « tte » = toute; « sss » = si et seulement si; his « ⟷ »
% between τ_i and τ_j (p. 2) means « commute with ». His boxed category names
% (C_0, C'_0, C''_0, C'''_0) are set \boxed. Two checks of his algebra are
% flagged in notes (p. 17 signs of the cross terms of f_V; p. 18 c') written
% « = id » where the gloss requires « ≠ id »); p. 19's β_0 + β_1 - 4 =
% α_0 + α_1 checks.
%
% The dating is the folder's own; the group « Géométrie et topologie
% combinatoire » (69-89) holds it.
\documentclass[11pt,a4paper]{article}
\input{../preamble/grothendieck}

\folder{72}
\batch{1}
\pages{1}{20}
\dating{[à partir de 1978]}
\watermark{Édition de démonstration}
\foldertitle{Polyèdres réguliers et hyperpolyèdres réguliers (Calculs en dimension quelconque) : notes manuscrites (s.d.).}

\begin{document}

\section{$n$-polyèdres réguliers et représentations de $\mathfrak{G}_n$ dans $O(n+1)$ \ldots\ (dictionnaire)}
\note{titre de sa main, en haut à droite de la feuille de chemise qui est la
page 1 ; le signe de « $n+1$ » est d'une lecture incertaine (peut-être
$n-1$). À gauche, au crayon, « (23 p) »}

\page{2}
$n \geq 1$ fixé, et $n$-hyperpolyèdre combinatoire épinglé universel $\Pi$,
groupe d'autom.
\[
\mathfrak{G}_n = \lbrace \tau_0, \ldots, \tau_n \mid \tau_i \leftrightarrow \tau_j
\ \text{si}\ 0 \leq i, j \leq n,\ i, j\ \text{non cons.} \rbrace
\]
\note{la flèche double $\tau_i \leftrightarrow \tau_j$ note ici, comme dans la
condition a) plus bas, que $\tau_i$ et $\tau_j$ commutent}

$S$ schéma de base fixé

\marginal{autodual} $C_{\mathrm{pr}}$ cat.\ des réal.\ géom.\ régulières
\add{non dég.} projectives de $\Pi$ (autodual)

\marginal{autodual} $C'_{\mathrm{pr}}$ cat des systèmes \struck{\ill{}}
$(X, (\tau_i)_{0 \leq i \leq n},$ \struck{\ill{}} $X_{*})$
\add{\uncertain{$\tau_j \mapsto \pm 1$}}
\note{l'ajout est écrit sous la ligne, sous le passage biffé, où se lisait
peut-être $\varphi_G : \mathfrak{G}_n \to \ldots$}

$X$ fibré projectif, $\tau_i$ $(0 \leq i \leq n)$ \uncertain{ps.}\ réfl.\ strictes
$(\Rightarrow h_i, H_i)$ \struck{, $s_0$ section}

$X_{*}$ drapeau maximal, avec conditions
\begin{itemize}
\item[a)] $\tau_i \leftrightarrow \tau_j$ si $0 \leq i, j \leq n$ \quad $i, j$
non consécutifs.
\item[b)] $\tau_i X_j = X_j$ si $0 \leq i \neq j \leq n$
\item[c)] $\tau_i X_i \neq X_i$ \quad $0 \leq i \leq n$
\end{itemize}

\marginal{NB $C_{\mathrm{pr}}$, $C'_{\mathrm{pr}}$ proviennent par
\emph{rigidification} de catégories $\widetilde{C}_{\mathrm{pr}}$ \struck{\ill{}}
de \emph{représentations} projectives de $\mathfrak{G}_n$ sans plus}
\note{note de marge entourée, en regard des conditions a) à c)}

\marginal{NB Au lieu de se donner $X_{*}$, il suffit de se donner \struck{un}
$X_i$ avec $\tau_j X_i = X_i$ $(j \neq i)$ $\tau_i X_i \neq$ \uncertain{$0$}
satisfaisant à conditions \ill{} — p.\ ex.\ $s_0 = X_0$, tel que
\struck{\ill{}} \add{$(s_0, (h_i))$} \struck{\ill{}} engendre linéairement
$X$}

\marginal{non autodual} $C_{\mathrm{aff}}$ variante affine de $C_{\mathrm{pr}}$
[NB. il n'est plus besoin de se donner $\varphi_G$ \ill{} pour
$C'_{\mathrm{aff}}$]

\marginal{non autodual} $C'_{\mathrm{aff}}$ variante affine de
$C'_{\mathrm{pr}}$

Attention, $\tau_0$ \uncertain{n'est peut-être pas} une ps.\ réflexion affine,
mais $\widehat{\tau}_0$ est une ps.\ refl.\ projective (dont l'hyperplan peut
être exceptionnellement l'hyperplan à l'$\infty$ \ldots)
\marginal{car $\beta_0 = \mu_0 = 0$ ! (\ill{} en \uncertain{car.\ 2}..)}

\marginal{autodual} $C_{\mathrm{vect}}$ variante vectorielle de
$C_{\mathrm{pr}}$ (où $X$ \uncertain{rigidifié} comme $P(\mathcal{E})$, et $s_0$
rigidifié par une section \emph{de} $\mathcal{E}$ — ce qui
\struck{rigidifie} est une donnée de \emph{rigidification} dans une donnée de
catégorie $\widetilde{C}_{\mathrm{vect}}$, éliminant le groupe $G_m$
d'autom..)
\marginal{$\widetilde{C}_{\mathrm{vect}}$, $\widetilde{C}'_{\mathrm{vect}}$
sans \uncertain{$G_m$-épinglage}}

$C'_{\mathrm{vect}}$ \struck{variante} variante vectorielle de
$C'_{\mathrm{pr}}$ (on \uncertain{laisse tomber} \ill{} = le drapeau,
\uncertain{remplacé par} $h_0 = a_{-1}$ ---)

\struck{\ill{}} $C^{\infty}_{?}$ cat des
$(C, \varphi^{C}_{G} : \mathfrak{G}_n \to \operatorname{Aut} C, h_0)$ avec $C$
fibré projectif de \uncertain{rg $n$}, $\varphi^{C}_{G}$ hom groupes, $h_0$
section de $C$, avec
\note{l'indice de $C^{\infty}$, rendu ici par un point d'interrogation, est illisible ; les pages suivantes ont
$C^{\infty}_{\mathrm{vect}}$ et $C^{\infty}_{\mathrm{pr.orth}}$}
\begin{itemize}
\item[a)] les $\tau^{C}_{i}$ ps.\ réflex.\ $(0 \leq i \leq n)$ strictes si
$1 \leq i \leq n$, \uncertain{associée} à $h_0$ pour $\tau^{C}_{0}$ (qui peut
être non strict) $(\Rightarrow h_i, H_i$ pour $1 \leq i \leq n$ — mais pas de
$h_0)$
\end{itemize}
\marginal{non autodual, sauf si \ill{}, i.e.\ $c \notin C$
\uncertain{en tte fibre}}
\marginal{NB La donnée de $s_0$ superflue si $\tau^{C}_{0}$ stricte i.e.\
$\mu_0, \beta_0$ non nuls simultanément}

\page{3}
\begin{itemize}
\item[b)] $h_i$ $(0 \leq i \leq n)$ engendrent $C$ i.e.\ lin.\ indép.
\item[c)] $h_i \notin H_{i+1}$ $(0 \leq i \leq n-1)$ en tte fibre
\end{itemize}

$\widetilde{C}^{\infty}_{\mathrm{vect}}$ variante vectorielle
$(V, \varphi^{V}_{\mathfrak{G}}, L_0)$, $V$ fibré vect.\ rg $n+1$,
$\varphi^{V}_{\mathfrak{G}} : \mathfrak{G} \to \operatorname{Aut}(V)$, $L_0$
sous-fibré vect.\ de rg 1, avec
\marginal{non autodual sauf si $\rho$ inv.}
\begin{itemize}
\item[a)] les $\tau^{V}_{i}$ $0 \leq i \leq n$ pseudoréflexions strictes pour
$1 \leq i \leq n$, et ass.\ à $L_0$ si $i = 0$ ($\Rightarrow L_i, H_i$,
\uncertain{$L'_i$} si $1 \leq i \leq n$)
\item[b)] $V = \bigoplus_{0 \leq i \leq n} L_i$
\item[c)] $L_i \not\subset H_{i+1}$ \quad $(0 \leq i \leq n-1)$
\end{itemize}

Alors
$L_0 \simeq L_1^{\vee} \simeq L_1 \simeq L_2^{\vee} \simeq L_2 \cdots \simeq
L_n^{\vee} \simeq L_n$, soit $L$. On \uncertain{appelle}
\emph{rigidification} la donnée d'une base de $L$, (p.\ ex.\ d'une base
\ill{} de l'un des $L_i$, \uncertain{ou des $L'_i$})

$C^{\infty}_{\mathrm{vect}}$ cat.\ des objets de
$\widetilde{C}^{\infty}_{\mathrm{vect}}$ \emph{rigidifiés}

$C^{\infty}_{\mathrm{pr.orth}}$ cat.\ des objets
$(C, Q, \varphi^{C,Q}_{\mathfrak{G}}, h_0)$, avec
\marginal{$n \geq 2$, les \uncertain{$\mu_i = 2$}}
\note{note de marge entourée}
\marginal{non autodual sauf si $\rho$ inv.}
\struck{$C$, $\varphi^{C}_{\mathfrak{G}}$ objet de \ill{}, $Q$ hyperquadrique
\ill{} de $C$}
$C$, fibré projectif rg $n$, $Q$ \uncertain{hyperquadrique} \ill{},
$\varphi^{C,Q}_{\mathfrak{G}} : \mathfrak{G}_n \to \operatorname{Aut}(C, Q)$,
$h_0$ section de $C$, tels que
\begin{itemize}
\item[a)] les $\tau^{C,Q}_{i}$ $(0 \leq i \leq n)$ réfl.\ orth., strictes pour
$1 \leq i \leq n$, ass.\ à $h_0$ si $i = 0$ $(\Rightarrow h_i, H_i$
$1 \leq i \leq n)$
\item[b)] $h_i$ $(0 \leq i \leq n)$ engendrent $C$
\item[c)] $h_i \notin H_{i+1}$ en tte fibre $(0 \leq i \leq n-1)$
\end{itemize}

\marginal{NB Dans le cas \struck{\ill{}} \add{$\beta_0 \beta_1$ inv./plus}
(\struck{\ill{}} \add{\ill{}} inv., et
\struck{$(\beta_1, 2)\mathcal{O}_S = \mathcal{O}_S$}
\struck{($\beta_1$ inv.\ \ill{})}
$\Rightarrow$ $(\beta_1, 4 - \beta_0)\mathcal{O}_S = \mathcal{O}_S$ i.e.\
$\beta_1 = \alpha_1 + 2$ et $4 - \beta_0 = 2 - \alpha_0$ non nuls
simultanément) b) équivaut au fait que, \uncertain{sur chaque fibre}, il
\ill{} d'hyperplan \ill{} \struck{\ill{}} par $\mathfrak{G}$ \ldots\ (quelle
remarque valant aussi pour $C^{\infty}_{\mathrm{vect}}$,
$C^{\infty}_{\mathrm{pr}}$ \ldots\ explicitation}
\note{longue note de marge au bas de la page, séparée du texte par un trait
oblique ; plusieurs mots en sont biffés ou surchargés}

\emph{RemNB} si $\tau_i$, $\tau_{i+1}$ réflexions \emph{strict.}\ orth.\
i.e.\ \ill{} inv. Alors $h_i \notin H_{i+1}$ en tte fibre équivaut à
$\tau_i \tau_{i+1} \neq \tau_{i+1} \tau_i$ \add{en tte fibre} i.e.\
$\rho_i^{2} \neq \mathrm{id}$ en tte fibre. Si donc $\rho_i$ inv.\ (c) se
remplace par ces conditions pour $0 \leq i \leq n-1$ \ldots

\page{4}
\emph{RemNB} On \uncertain{aura même} $h_0 \in \Gamma(C - Q)$ \ldots

\marginal{non autodual sauf si $\rho$ inv.}
$C^{\infty}_{\mathrm{vect.orth}}$ cat.\ des syst.\
$(V, f, \varphi^{V}_{\mathfrak{G}}, a_0)$, avec $V$ fibré vectoriel de rg
$n+1$, $f$ forme quadratique \struck{non dég.\ en chaque fibre},
$\varphi^{V\,\mathrm{orth}}_{\mathfrak{G}} : \mathfrak{G} \to
\operatorname{Aut}(V, f)$, $a_0$ section de $V$, satisfaisant
\begin{itemize}
\item[a)] les $\tau^{V}_{i}$ \add{$(0 \leq i \leq n)$} réflexions orth., les
$\tau_i$ $(1 \leq i \leq n)$ \struck{non} $\neq \mathrm{id}$ en chaque fibre
$(\Rightarrow L_i, H_i$ pour $1 \leq i \leq n)$, $\tau_0$ associé à
$\mathcal{O}_S a_0$, \struck{\ill{}} et $\boxed{f(a_0) = 1}$ (de sorte que
$\mathcal{O}_S a_0$ \add{\ill{}} \uncertain{sous-fibré vect.}\ldots)
\item[b)], c) comme dessus, avec commentaires correspondants.
\end{itemize}

On pourrait définir $\widetilde{C}^{\infty}_{\mathrm{vect.orth}}$, catégorie
des \struck{systèmes} $(V, f, \varphi^{V}_{\mathfrak{G}}, L_0)$, avec
\struck{\ill{}} \struck{\ill{} comme dans $C^{\infty}_{\mathrm{vect}}$,}
$(V, f)$ \struck{\ill{}} \uncertain{fibré quadratique} de rg $n+1$,
$\varphi^{V}_{\mathfrak{G}} : \mathfrak{G}_n \to \operatorname{Aut}(V, f)$,
$L_0$ s-fibré de rg 1, tels que a) b) c) comme précédemment, mais
$f(a_0) = 1$ remplacé par $f \vert L_0$ \ill{} \emph{\uncertain{lisse}}.

Et \ill{} trouve $C^{\infty}_{\mathrm{vect.orth}}$ à partir de
$\widetilde{C}^{\infty}_{\mathrm{vect.orth}}$ par rigidification par une
section $a_0$ de $L_0$ telle que $f(a_0) = 1$ (les sections formant un
torseur sous $\mu_2$ \ldots).

\[
C_{\mathrm{pr}} \approxeq C'_{\mathrm{pr}} \approxeq C_{\mathrm{aff}}
\approxeq C'_{\mathrm{aff}} \approxeq C_{\mathrm{vect}} \approxeq
C'_{\mathrm{vect}} \;\Big\vert\; \approxeq C^{\infty}_{\mathrm{pr}}
\]
\[
\approxeq C^{\infty}_{\mathrm{vect}} \approxeq \;\Big\vert\;
C^{\infty}_{\mathrm{pr.orth}} \approxeq C^{\infty}_{\mathrm{vect.orth}}
\]
\note{les signes $\approxeq$ rendent un $\approx$ posé sur un long trait
horizontal. Deux traits verticaux coupent la chaîne ; au premier est accrochée
la mention « $n \geq 2$ (explicités variante si $n = 1$ !) », au second « si
de plus les $\lambda_i = -1$ i.e.\ les $\tau_i$ des réflexions. »}

\page{5}
Soit $G_0$ l'image de $\mathfrak{G}_n$ dans
$\operatorname{Aut}(\mathcal{E}, V)$ (un mot biffé, \ill{}, avant $V$), on a un diagramme

\begin{tikzcd}[column sep=small]
 & \operatorname{Aut}(X, C) \arrow[r] & \operatorname{Aut}(C) \\
G_0 \arrow[r, hook] & \operatorname{Aut}(\mathcal{E}, V, f) \arrow[u] \arrow[r] & \operatorname{Aut}(V, f) \arrow[u]
\end{tikzcd}

\note{au-dessus de $\operatorname{Aut}(X, C)$ est écrit
$\operatorname{Aut}(\mathcal{E}, f, \pi, \pi')$ ; en tête de la ligne du haut
un « $G_0 \to$ » est biffé, et plusieurs lettres (un indice sous
$\operatorname{Aut}(X, C)$, les seconds arguments des deux
$\operatorname{Aut}$ de droite) sont surchargées ou biffées}

On cherche \struck{\ill{}} \uncertain{les noyaux} de
\[
G_0 \longrightarrow
\left\lbrace
\begin{array}{l}
\operatorname{Aut}(X, \ldots) \\
\operatorname{Aut}(V)
\end{array}
\right.
\longrightarrow \operatorname{Aut}(C)
\]
\note{le second argument du premier $\operatorname{Aut}$ est biffé}

a) $\operatorname{Ker}(\operatorname{Aut}(\mathcal{E}, V, f) \to
\operatorname{Aut}(X))$ est formé des homothéties qui conservent $f$, donc
isom.\ $\mu_2$. En fait, $G_0$ conserve mieux que $(\mathcal{E}, f, V)$, il
conserve $(\mathcal{E}, f, \pi, \pi')$ \add{\ill{} $f$}, donc si c'est une
homothétie c'est l'identité : $G_0$ opère \emph{fidèlement} sur $X$.

b) $\operatorname{Ker}(\operatorname{Aut}(\mathcal{E}, f, \pi, \pi') \to
\operatorname{Aut}(V))$ est formé des \ill{} \ill{} de la forme
\[
\mathrm{id} + \pi' \otimes a \qquad (a \in \Gamma \mathcal{E}),
\]
\note{devant $\pi'$, un signe biffé, illisible}
qui conservent $\pi, \pi', f$. Ils conservent $\pi$ sss \struck{\ill{}}
$\rho a = 0$ (donc $a = 0$ si $\rho$ inv., plus gén.\ si $\rho$ est régulier)
(NB $\langle \pi, \pi' \rangle = \rho$), et ils conservent $\pi'$ sss :
\[
\langle a, \pi' \rangle \pi' = \ldots
\quad \text{i.e.} \quad \langle a, \pi' \rangle = 0
\quad \text{i.e.} \quad \boxed{a \in V.}
\]
\note{trois courts passages biffés et illisibles dans cette ligne : après
le premier $\langle a, \pi' \rangle$, au second membre de la première
égalité (rendu ici par $\ldots$), et devant le second
$\langle a, \pi' \rangle$}
Conserver $f$ signifie
\[
f(a)\pi' + \psi(a) = 0
\]
\note{la formule porte à gauche un repère entouré, peut-être « 8 »}
d'où \struck{\ill{}} en \ill{} $\psi'$
\[
f(a) \underbrace{\psi'(\pi')}_{-\sigma' \pi} +
\underbrace{\psi' \psi(a)}_{\beta a} = 0
\quad \text{i.e.} \quad \beta a = \sigma' f(a) \pi
\]
\note{devant $\sigma'$ dans le dernier membre, une petite marque illisible,
peut-être biffée}

\page{6}
Si $\beta$ inversible ça conduit $a = \lambda \pi$, dans l'équation
précédente \uncertain{(8)} devient
\[
f(a)\pi' + \lambda \underbrace{\psi(\pi)}_{-\sigma \pi'} = 0
\quad \text{i.e.} \quad
\sigma \lambda = \lambda^{2} \underbrace{f(\pi)}_{-\widetilde{\sigma}\widetilde{\rho}}
\]
\struck{$\beta \lambda \pi = \sigma' \lambda^{2} f(\pi) \ldots$}

i.e.
\[
\beta \lambda = \sigma' \lambda^{2} \underbrace{f(\pi)}_{-\widetilde{\sigma}\widetilde{\rho}}
\qquad \text{i.e.} \quad \lambda(\sigma + \lambda \widetilde{\sigma}\widetilde{\rho}) = 0
\]
\note{les deux colonnes sont écrites côte à côte : à gauche la variante en
$\beta$, $\sigma'$, à droite celle en $\sigma$}
i.e.
\[
\lambda(\beta + \sigma' \widetilde{\sigma} \widetilde{\rho} \lambda) = 0,
\quad \text{en plus de } \rho \lambda = 0
\]
Si on est dans un anneau intègre, cela veut que \emph{soit} $\lambda = 0$
i.e.\ $u = \mathrm{id}$, soit $\beta + \sigma' \widetilde{\sigma}
\widetilde{\rho} \lambda = 0$ ce qui implique (comme $\beta$ inv.\ par hyp.)
$\sigma' \widetilde{\sigma} \widetilde{\rho}$ inversible (et
$\lambda = -\frac{\beta}{\sigma' \widetilde{\sigma} \widetilde{\rho}}$)
absurde si $\rho = \rho$ \ill{} (on impose) qui implique $\lambda = 0$
(puisque $\rho \lambda = 0$), \uncertain{comme}
$\rho = 2\widetilde{\rho}$, si $2$ inv.\ \struck{\ill{}} \add{\ill{}} ce
résultat. Mais si \ill{} \ill{}, \struck{\ill{}} $\rho = 0$ dans
$\rho \lambda = 0$ pour en conclure, il reste la condition (\struck{\ill{}} et
\ill{}) $\sigma + \lambda \widetilde{\sigma} \widetilde{\rho} = 0$ \ill{} on
$\sigma \sigma' = \beta + \gamma \rho = \beta$ \struck{\ill{}} \ill{}, donc
\add{$\sigma$, $\sigma'$ et} $\widetilde{\sigma} = \widetilde{\sigma}'$ \ill{}
inv., la condition devient $1 + \lambda \widetilde{\rho} = 0$ i.e.\
$\widetilde{\rho}$ inv.\ et $\lambda = -\frac{1}{\widetilde{\rho}}$. Ainsi
\note{« absurde si $\rho = \rho$ » : le second membre porte peut-être un
accent ; lecture incertaine}

\emph{Prop} Sur un anneau intègre, si $\beta$ ou $\rho$ \struck{$\neq 0$},
alors le seul $u \in \operatorname{Aut}(\mathcal{E}, f, \pi, \pi')$ induisant
l'identité dans $V$ est l'identité, \emph{sauf} dans le cas \struck{\ill{}}
$\beta \neq 0$, \uncertain{$n$ pair}, \struck{$(\rho = 2\rho')$} car.\ 2
$(\rho = 2\rho' = 0)$, $\rho'$ inv.\ auquel cas le seul autre $u$ est
\[
\mathrm{id} - \frac{1}{\widetilde{\rho}}\, \pi' \otimes \pi
= \mathrm{id} + \frac{1}{\widetilde{\rho}} \langle \pi', \ldots \rangle
\]
i.e.\ l'antipodisme
\note{le premier signe est surchargé d'une tache, le second pourrait être un
$\pm$ ; les arguments du crochet, au bord de la feuille, sont illisibles.
L'énoncé écrit $\rho'$ là où la formule écrit $\widetilde{\rho}$}

\page{7}
[Il reste à voir ce que ça donne si $\beta = 0$, $\rho = 0$ donc
$\sigma \sigma' = 0$. Posons
\[
a = \sum_{0}^{n} \lambda_i a_i, \quad \text{d'où} \quad
\psi(a) = \sum -\lambda_i (\beta_0 - \beta_{i-1}) a_i,
\]
\[
f(a) = \sum_{0}^{n} \lambda_i^{2} (\beta_0 - \beta_{i-1}) + (\lambda_0
\lambda_1 \beta_0 + \lambda_1 \lambda_2 (\ldots) + \cdots + \lambda_{n-1}
\lambda_n (\beta_0 - \beta_{n-1}))
\]
\note{les indices des $\beta$ sont d'une lecture incertaine ($\beta_{i-1}$ ou
$\beta_{i+1}$) ; le coefficient de $\lambda_1 \lambda_2$ est surchargé et
illisible, et le signe qui ouvre la seconde somme est repassé}
———— ]

Regardons enfin le noyau de
\[
\operatorname{Aut}(\mathcal{E}, f, \pi, \pi') \longrightarrow
\operatorname{Aut}(\overset{C}{\overbrace{P(V)}})
\]
\note{le $C$ est écrit au-dessus de $P(V)$, qui est donc $C$}
\ill{} les $u$ \ill{} \ill{} qui induisent une homothétie sur $V$. Cette
homothétie $\varepsilon$, conservant $f \vert V$, doit être dans $\mu_2$,
donc $\varepsilon u$ \struck{conserve} est l'identité sur $V$, donc de la forme
\[
\varepsilon u = \mathrm{id}_{\mathcal{E}} +
\pi' \otimes \varepsilon a \qquad (a \in \mathcal{E})
\]
\note{un $\varepsilon$ biffé devant $\pi'$}

\struck{A ——— $\varepsilon\,\mathrm{id}$.}

\struck{$u =$} $u = \varepsilon\,\mathrm{id} + \pi' \otimes a$ (un $\varepsilon$ biffé devant $a$)
\quad conserve $f, \pi, \pi'$.

Conserve $\pi, \pi'$
\[
\pi = \varepsilon \pi + \rho a \quad \text{i.e.} \quad
\boxed{(1 - \varepsilon)\pi = \rho a}
\]
\[
\pi' = \varepsilon \pi' + \langle a, \pi' \rangle \pi'
\quad \text{i.e.} \quad
\boxed{\langle a, \pi' \rangle = 1 - \varepsilon}
\]
\note{dans chacune des deux lignes, un $\varepsilon$ biffé devant le
dernier terme du second membre. À droite de la seconde ligne, un premier cadre
$\varepsilon(1 + \langle a, \pi' \rangle) = 1$ est biffé ; dans le cadre
retenu, un $\varepsilon a$ devant le crochet est biffé, et le $1$ du second
membre est repassé}

Enfin, conservation de $f$ \add{\ill{}} (signifie)
\[
\boxed{f(a)\pi' + \varepsilon \psi(a) = 0}
\]

\page{8}
\note{la moitié supérieure de la page, jusqu'au calcul entouré inclus, est
barrée de longs traits obliques ; elle se lit et est donnée ici}
Plaçons-nous sur un corps $k$, et supposons $\varepsilon \neq 1$. (donc car.\
$\neq 2$) — le cas $\varepsilon = 1$ étant le cas déjà traité (sous la
restriction que $\beta, \rho$ non nuls tous deux). On doit donc avoir
\emph{$\rho$ inv.}\ par la première équation, donc
\[
a = \frac{\ldots}{\rho}\, \pi = \lambda \pi
\quad \Bigl(\lambda = \frac{2}{\rho}\Bigr)
\]
\note{le numérateur de la première fraction, biffé, se lit
$1 - \varepsilon$ ; un mot biffé et illisible suit la parenthèse}
la 2\textsuperscript{e} équation donne $\rho \lambda = 2$ déjà satisfaite,
enfin la dernière s'écrit
\[
\lambda^{2} \underbrace{f(\pi)}_{-\widetilde{\rho}\widetilde{\sigma}} \pi'
+ \varepsilon \lambda \underbrace{\psi(\pi)}_{-\sigma \pi'} = 0
\]
i.e.
\[
\lambda(\underbrace{\widetilde{\rho}\widetilde{\sigma}\lambda}_{\frac{2\widetilde{\rho}\widetilde{\sigma}}{\rho} = \frac{\rho\widetilde{\sigma}}{\rho} = \widetilde{\sigma}} + \varepsilon \sigma) = 0
\]
\note{le dernier membre de l'accolade, $\widetilde{\sigma}$ ou $\sigma$, est
d'une lecture incertaine}

\struck{ce qui exige \ill{} $\lambda \neq 0$ \ill{}}
\[
\widetilde{\rho}\widetilde{\sigma}\lambda + \varepsilon \sigma = 0
\]
\note{dans ces trois formules, le signe $+$ devant $\varepsilon$ est biffé
sur la page, sans rien qui le remplace ; il est donné tel qu'écrit ; dans la deuxième, un signe
biffé et illisible ouvre la parenthèse}
\uncertain{ou encore} \struck{\ill{}} $\sigma = \widetilde{\rho}\widetilde{\sigma}$
\quad \struck{$(\ldots)\frac{2\widetilde{\rho}\widetilde{\sigma}}{\rho} + \varepsilon\sigma = 0$}
\note{ces trois lignes sont enfermées dans une grande boucle et en partie
griffonnées}

Sur un anneau intègre, les seuls cas pour $\varepsilon$ sont $\varepsilon = +1$
(déjà étudié, sous la restriction que $\rho, \beta$ non tous deux inv.) et
$\varepsilon = -1, \neq +1$ (donc \struck{car.\ $\neq 2$}) auquel cas \ill{}
\ill{} de car.\ $\neq 2$) d'où
\[
\rho a = 2\pi \quad \text{donc} \quad a = \frac{2}{\rho}\pi, \qquad
\varepsilon u = -\mathrm{id} + \pi' \otimes \frac{2}{\rho}\pi
\]
\note{sous $a = \frac{2}{\rho}\pi$, une mention de deux mots, illisible,
reliée par un trait à la ligne suivante ; dans $\varepsilon u$, de petites
marques en exposant devant $\pi'$ et devant la fraction, illisibles}
c'est l'antipodisme, dont on sait bien qu'il conserve $f$ et \struck{\ill{}}
$\pi, \pi'$ \ldots\ Donc

\struck{Proposition} \add{(Corps de base)} \struck{Prop} Si $\beta$ ou $\rho$
est inv.\ \struck{\ill{}} alors le seul automorphisme de
$(\mathcal{E}, f, \pi, \pi')$ qui induise l'identité

\page{9}
sur $C = \check{P}(V)$, est l'\struck{antiidentité} identité et l'antipodisme
(défini \struck{\ill{}} si $\widetilde{\rho}$ inv.\ \struck{(car.\ \ill{}
\ill{} pair)} ou si $\rho$ \ill{} comme
\[
\underline{a} = \mathrm{id}_{\mathcal{E}} - \frac{1}{\widetilde{\rho}}\,
\pi' \otimes \pi \ \text{si $n$ pair}, \qquad
\underline{a} = \mathrm{id}_{\mathcal{E}} - \frac{2}{\underset{\widetilde{\rho} = \rho}{\widetilde{\rho}}}\,
\pi' \otimes \pi \ \text{si}
\]
$n$ impair — dans ce dernier cas, c'est d'ailleurs l'identité en car.\ 2.]
\note{le $C$ initial est repassé ; le crochet fermant répond peut-être à celui
qui ouvre la page 7. Le reste de la page est blanc}

\page{10}
\emph{Dualité} \struck{\ill{}}

La condition $\rho$ \uncertain{inversible} suffit pour que l'accouplement
naturel entre $V$ et $V'$ soit une dualité — elle implique donc que $C$ et
$C'$ sont en dualité l'un avec l'autre. Si de plus $\beta$ inv.\ i.e.\
\struck{\ill{}} $f \vert V$ non dég.\ i.e.\ $\psi_V : V \simeq V'$, on trouve
un iso entre \struck{\ill{}} $C$ et $C'$ qui commute : $\mathfrak{G}$ et
transforme donc $Q_C$ en $Q_{C'}$. Donc dans ce cas
\[
(C, Q_C) \simeq (C', Q_{C'})
\]

Le passage d'un \struck{\ill{}} polyèdre régulier projectif épinglé, avec
son dual, sous ces conditions, en termes de $C^{\infty}_{\mathrm{pr.orth}}$,
s'exprime donc par la conservation de $(C, Q_C)$, et en
\struck{\ill{}} \add{remplaçant} $\tau^{C}_{0} \ldots \tau^{C}_{n}$ par
$\tau^{C}_{n} \ldots \tau^{C}_{0}$ (i.e.\
$\varphi^{C}_{\mathfrak{G}}$ par $\varphi^{C}_{\mathfrak{G}} \circ \theta$, où
$\theta : \mathfrak{G} \to \mathfrak{G}$ \uncertain{échange} $\tau_i$ et
$\tau_{n-i}$) — ici il est inutile de préciser une section $a_0$
\add{\uncertain{rigide}} de $C$, puisque $\beta_0, \beta_n$ inversibles.
\note{le reste de la page est blanc}

\page{11}
Soit maintenant $\Pi$ hyperpolyèdre comb.\ de dim.\ $n$, $G$ le groupe de
ses automorphismes, $R_{\Pi}$ \struck{$G$}-torseur des repères — son
\uncertain{commutant} est un quotient \add{\uncertain{privilégié}}
$\mathfrak{G}_{\Pi}$ de $\mathfrak{G}$,

$C_{\mathrm{pr}}(\Pi)$ cat.\ des réal.\ géom.\ régulières non dégénérées de
$\Pi$ $\Longrightarrow \varphi_G : G \to \operatorname{Aut}(X, \ldots)$.
\note{la flèche $\Longrightarrow$ et la suite sont écrites sous la ligne et
reliées par un trait à « non dégénérées »}

\struck{\ill{}}

À tt $r \in R$ correspond un objet de $C_{\mathrm{pr}}$ (\uncertain{du}
universel épinglé) et pour $r \in R$ variable, on a un système transitif
d'iso.\ entre ces objets. Donc $C_{\mathrm{pr}}(\Pi)$ s'identifie
\add{(pour $\Pi$ fixé)} à la \add{sous-}cat.\ pleine \struck{des couples
$(\ldots)$} de $C_{\mathrm{pr}}$ formée des $(X, \varphi_{\mathfrak{G}},
\ldots)$ tels que $\varphi_{\mathfrak{G}} : \mathfrak{G} \to
\operatorname{Aut}(X)$ se factorise par le quotient $\mathfrak{G}_{\Pi}$ de
$\mathfrak{G}$. On fait opérer \emph{à droite} $\mathfrak{G}_{\Pi}$ sur
$R_{\Pi}$ via son action à gauche, et on \uncertain{s'établit} par
\uncertain{tordre} $(X, \varphi_{\mathfrak{G}}, \ldots)$ en
\[
R_{\Pi} \wedge^{\mathfrak{G}_{\Pi}} X
\]
Même procédé pour les autres descriptions — par $C'_{\mathrm{pr}}(\Pi)$
etc.\ \ldots

Autre point de vue : On ne se donne pas $\Pi$, mais seulement un
\add{groupe} quotient \uncertain{déterminé} $\mathfrak{G}_0$ de
$\mathfrak{G}$, \struck{\ill{}} correspondant à un hyperpolyèdre régulier
combinatoire épinglé. On considère ses hyperpolyèdres réguliers

\page{12}
combin.\ de type $\mathfrak{G}_0$ sur $S$ comme la donnée d'un torseur
\add{$R_0$} sous $\mathfrak{G}_S$, permettant de tordre $\Pi_0$. Une
représentation de ce \struck{\ill{}} $\Pi$ \add{\struck{\ill{}}} (variant :
la donnée (en plus de $\Pi$) d'un objet de $C_{\mathrm{pr}}$ ou d'une des
cat.\ équivalentes. Donc la cat.\ des réal.\ géom.\ projectives (disons)
\struck{dans} $C$ des hyperpolyèdres comb.\ réguliers tordus sur $S$, de type
$\mathfrak{G}_0$, est équivalente à la cat.\ produit de
$C_{\mathrm{pr}}$\struck{\ill{}} (ou d'une de ses variantes) et de la cat.\
des $\mathfrak{G}_0$-torseurs.
\note{la parenthèse ouverte après « variant » n'est pas refermée ; le
$\mathfrak{G}_S$ de la deuxième ligne est peut-être $\mathfrak{G}_{0}$ ou
$G_S$. Sous le texte, un long trait de crayon ondulé descend en diagonale ;
le reste de la page est blanc}

\page{13}
$\Pi$ \uncertain{$n$-}\uncertain{quasi-}polyèdre comb.\ \emph{épinglé} universel

$G = \mathfrak{G}_n$ son groupe d'autom., engendré par
$(\tau_i)_{0 \leq i \leq n}$

$S$ schéma de base

\marginal{NB $C$ catégorie rigide, dont les classes d'iso.\ sont en corr.\
1-1 avec $\Gamma(S, \mathcal{O}_S)^{2n-1}$ via $\lbrace \lambda_0, \alpha_0,$ $\lambda_1, \alpha_1,$ $\ldots,$
$\lambda_{n-1}, \alpha_{n-1},$ $\lambda_n \rbrace$}
\note{l'exposant $2n-1$ est d'une lecture incertaine ; la liste qu'il compte
a $2n+1$ termes}

$\boxed{C_0}$ = catégorie \add{(groupoïde)} des réalisations géom.\ non dég.\
projectives sur $S$ $\simeq$ Catégorie des triples
$(X,$ \struck{$(\tau_i)$} $(\tau_i)_{0 \leq i \leq n}, s_0)$, $X$ fibré
projectif de rg $n+1$, sur $S$, $\tau_i$ une pseudo-réflexion stricte de $X$
(d'où centre $h_i$, co-centre $H_i$) $0 \leq i \leq n$, $s_0$ section de $X$
sur $S$, avec conditions
\begin{itemize}
\item[1)] \struck{\ill{}} $(h_i)_{0 \leq i \leq n}$ \struck{engendrent}
\add{lin.\ indép.}\ en chaque fibre ($\Longrightarrow$ $C$ hyperplan dans $X$)
\item[1')] \struck{\ill{}} $(H_i)$ lin.\ indép.\ en chaque fibre
($\Longrightarrow$ $c$, section de $X$)
\item[2)] $\tau_i \tau_j = \tau_j \tau_i$ si $0 \leq i, j \leq n$, $i, j$ non
consécutifs ($\Longleftrightarrow$ $h_i \subset H_j$ \struck{\ill{}} \& $0 \leq
i, j \leq n$, $i, j$ non consécutifs) \add{conséquence de 1, 5}
\item[3)] $h_i \notin H_{i+1}$ en \struck{\ill{}} tte fibre \quad
$0 \leq i \leq n-1$ \quad ($\Longrightarrow$ \struck{\ill{}} $\tau_i, \tau_{i+1}$
ne commutent pas)
\item[4)] $s_0 \subset H_i$ \quad $(1 \leq i \leq n)$
\item[5)] $s_0 \notin H_0$ en tte fibre \struck{\ill{}}
\end{itemize}
\note{« conséquence de 1, 5 » est écrit sous le $\Longleftrightarrow$ de la
condition 2) et lui est relié par un trait ; les chiffres sont d'une lecture
incertaine}

\marginal{NB 1') est conséquence des autres conditions 1) 2) 3) 4) 5) ; mais
en l'absence d'un $s_0$ ct 4) 5), il faut imposer 1') (ce qui revient
\uncertain{à 3)} \uncertain{rapportée} autrement)}
\marginal{\ill{} — une note de plusieurs lignes écrite en oblique entre la
marge et le texte, contre le trait vertical qui embrasse les conditions ; se
lisent « condition », « $s_0$ », « si $x$ », « \uncertain{il existe} »}

\marginal{NB Il y a lieu d'introduire aussi la cat.\ des couples
$(\mathcal{E}, \varphi^{\mathcal{E}}_{G})$ \uncertain{$(\mathcal{E},
\tau_i)$} (sans rig.\ $s_0$ \ldots)}

$\boxed{C''}$ = cat des triples $(V, \varphi^{V}_{G}, L_0)$, $V$ fibré
vectoriel de rg $n+1$ sur $S$, $\varphi^{V}_{G} : G \to \operatorname{Aut}(V)$
hom.\ de groupes, $L_0 \subset V$ sous-fibré inv., tels que
\note{le $V$ et le $\varphi$ du triple sont repassés sur d'autres lettres ;
le $L_0$ remplace un mot biffé}
\begin{itemize}
\item[1)] $\forall\, 0 \leq i \leq n$, $\varphi^{V}_{G}(\tau_i)$
\add{$\tau^{V}_{i} \overset{\mathrm{def}}{=}$} est une ps.\ réflexion de $V$,
stricte si $1 \leq i \leq n$, et \uncertain{associée} : $L_0$ si $i = 0$
(pas nécessairement stricte si $i = 0$). On en conclut des $L_i \subset V$,
$L'_i \subset V'$ (ss-fibrés inv.) $1 \leq i \leq n$
\item[2)] $V = \bigoplus L_i$, \struck{\ill{} [2') des $(L'_i)_{1 \leq i
\leq n}$ indép.]} en chaque fibre (i.e.\ $\bigcap_{1 \leq i \leq n} H_i$
\emph{droite} \uncertain{sur} chaque fibre] $\boxed{\ }$
\item[3)] \struck{$L_0 \subset H_i$ si $i \geq 2$ \quad (NB est conséquence
de 1), 2) si $\tau^{V}_{0}$ ps.\ refl.\ stricte)} \quad \struck{et \ill{}
\ill{} $H_i = (L'_i)^{\perp}$}
\end{itemize}
\note{le « $\Delta_d$ » écrit au-dessus de « droite » est d'une lecture
incertaine ; la condition 3) est barrée de traits obliques}

\marginal{NB 1), 2) $\Rightarrow L_i \subset H_j$ si $1 \leq i, j \leq n$
$i, j$ non cons.}
\marginal{NB 2') est conséquence des autres conditions \uncertain{y inclus}
3), et 3) est conséquence des autres \uncertain{y inclus}}

\page{14}
\begin{itemize}
\item[\struck{3}4)] $L_i \not\subset H_{i+1}$ \quad $0 \leq i \leq n-1$, en tte
fibre i.e.\ $L_i \otimes L'_{i+1} \xrightarrow{\ \sim\ } \mathcal{O}_S$
\end{itemize}
\note{sous la flèche, reliée à elle : « hom.\ can.\ induit par
$V \otimes \check{V} \to \mathcal{O}_S$ » (le premier facteur est repassé)}

\struck{$C'_0$}

\struck{$C'_0$ = Cat des triples précédents \emph{rigidifiés} par une section
$a_0$}

\emph{NB} De 4) on tire
\[
\left\lbrace
\begin{array}{ll}
L_i, L'_{i+1} \ \text{en dualité} & (0 \leq i \leq n-1) \ \text{donc} \\
L_i \xrightarrow{\ \sim\ } L'^{\vee}_{i+1} & (0 \leq i \leq n-1)
\end{array}
\right.
\]
De 1) on tire
\[
\begin{array}{l}
L'_i \otimes L_i \xleftarrow{\ \sim\ } \mathcal{O}_S \\
\tau^{V}_{i} - \mathrm{id}_V \longmapsto 1
\end{array}
\qquad \Big\vert \quad 1 \leq i \leq n
\]
\struck{\uncertain{dim.}} donc
\[
L'_i \underset{\mathrm{can}}{\simeq} L_i^{\vee} \qquad 1 \leq i \leq n
\]
On en conclut de proche en proche des isos
\[
L_0 \simeq L_1'^{\vee} \simeq L_1 \simeq L_2'^{\vee} \simeq L_2 \cdots
\simeq L_n'^{\vee} \simeq L_n
\]
d'où
\[
V \simeq \mathbf{L} \otimes \mathcal{O}_S^{\,n+1} \qquad \text{iso can.}
\]
\note{le $\mathbf{L}$ est écrit en gras sur la page et n'est pas défini ;
c'est le fibré en droites commun de la chaîne précédente}

$\boxed{C'_0}$ = cat des triples $(V, \varphi^{V}_{G}, L_0)$ \emph{rigidifiés}
par une section \add{-base} $a_0$ de $L_0$, i.e.\ des triples
$(V, \varphi^{V}_{G}, a_0)$, \struck{$a_0$ de} $V, \varphi^{V}_{G}$ comme
dessus, $a_0$ section $\neq 0$ en tte fibre, satisfaisant les conditions 1) à
4) où on fait $L_0 = \mathcal{O}_S a_0$ [le fait que $a_0$ est $\neq 0$ en
chaque fibre est déjà contenu dans \struck{\ill{}} la première condition 4),
si on l'écrit $a_0 \notin H_1$ \struck{en} tte fibre \ldots]

\page{15}
Un tel $V$ a donc une base canonique \add{$(a_0, a_1, \ldots, a_n)$}
(\uncertain{indépendante} par \uncertain{le} $a_0$ donné) de telle façon que,
\struck{si} $a_i$ forme base de $L_i$ \add{$0 \leq i \leq n$},
\struck{\ill{}} puis on ait
\[
\tau_i = \mathrm{id} + a_i^{\vee} \otimes a_i
\]
avec $a_i^{\vee} \in \Gamma \check{V}$ bien déterminé, et
\[
\langle a_i, a_{i+1}^{\vee} \rangle = 1 \qquad 0 \leq i \leq n-1
\]
(on a donc exprimé déjà \struck{1), 2)} sur $(\tau_i^{\vee})$ 1), 2)
\uncertain{3)}), il \struck{\uncertain{en reste deux}} reste : exprimer
\struck{\ill{}} la commutation des $\tau_i \tau_j$ pour $i, j$ non
consécutifs, ce qui s'exprime par
\[
L_i \subset \operatorname{Ker}(a_j^{\vee}) \quad \text{si} \quad
0 \leq i, j \leq n, \ i, j \ \text{non consécutifs.}
\]
On aura donc
\[
\left\lbrace
\begin{array}{l}
a_0^{\vee} = (\lambda_0 - 1) a_0^{*} + \beta_0 a_1^{*} \\[4pt]
a_i^{\vee} = \xi\, a_{i-1}^{*} + (\lambda_i - 1) a_i^{*} + \beta_i a_{i+1}^{*}
\qquad 1 \leq i \leq n-1 \\[4pt]
a_n^{\vee} = a_{n-1}^{*} + (\lambda_n - 1) a_n^{*}
\end{array}
\right.
\]
\note{les $a_i^{*}$ sont, semble-t-il, la base duale de $(a_i)$. Le premier
coefficient de la ligne médiane, rendu $\xi$, est un signe douteux ; la borne
$1 \leq i$ remplace un $0$ biffé ; la dernière ligne
porte $a_{n-1}^{*}$ sans coefficient. Dans la première ligne, un premier
membre $(\lambda_0 - 1) a_0^{*} + \beta_0 a_1^{*}$ est surchargé et récrit
au-dessus}

\marginal{NB $a_0^{\vee} = 0$ sss $\lambda_0 = 1$, $\beta_0 = 0$. Donc $\tau_0$
ps.\ réfl.\ stricte sss \struck{\ill{}} \add{\ill{}} $\forall\, s \in S$, on
a $\lambda_0(s) \neq 1$ ou $\beta_0(s) \neq 0$, i.e.\ sss
$(\lambda_0 - 1)\mathcal{O}_S + \beta_0 \mathcal{O}_S = \mathcal{O}_S$}
\note{note de marge entourée d'un trait}

\emph{Scholie} On a un foncteur $C \to C'_0$ évident
\[
(\mathcal{E}, \tau_{*}, s_0) \longmapsto (V_{*} \ldots
\]
qui est une équivalence de catégories
\note{le second membre de la flèche est inachevé}

$\boxed{C''_0}$ \struck{\ill{}} Considérons la catégorie des
\struck{\ill{}} \add{triples} $(C, \varphi^{C}_{G}, h_0)$, où $C$ est un
fibré projectif de rg $n$, $\varphi^{C}_{G} : G \to \operatorname{Aut}(C)$
hom.\ de groupes, $h_0$ section de $C$, avec conditions
\marginal{$n \geq 2$}
\begin{itemize}
\item[1)] $\forall\, 0 \leq i \leq n$, $\tau_i \overset{\mathrm{def}}{=}
\varphi^{C}_{G}(\tau_i)$ est une pseudo-réflexion, qui est stricte si
$1 \leq i \leq n$, et associée au centre $h_0$ si $i = 0$ (\struck{\ill{}} pas
nécess.\ stricte) — d'où $h_i, H_i$ $(1 \leq i \leq n)$
\item[2)] Les $h_i$ $0 \leq i \leq n$ lin.\ indép.\ en chaque fibre
\item[3)] $h_i \notin H_{i+1}$ en tte fibre \quad $(0 \leq i \leq n-1)$
\end{itemize}

\page{16}
On a un foncteur
\[
C'_0 \longrightarrow C''_0
\]
c'est une équivalence de catégories \emph{si $n \geq 2$} (mais non si
$n = 2$)
\note{ainsi sur la page : « si $n \geq 2$ » (souligné deux fois), puis « (mais
non si $n = 2$) » ; l'un des deux porte peut-être $n = 1$}

En effet, loc.\ on peut \uncertain{écrire} $C = \check{P}(V)$, et
remonter de façon unique \add{\uncertain{(si $n \geq 2$ !)}} les
\add{ps.\ réflex.}\ $\tau^{C}_{i}$ en des \add{ps.\ réfl.}\ $\tau^{V}_{i}$,
qui vont satisfaire aux \uncertain{mêmes} conditions de commutation. Le
module $V$ est déterminé modulo $\otimes L$, $L$ inv., et cette
indétermination est levée par la rigidification $L_0 \simeq \mathcal{O}_S$
\ldots
\note{« (si $n \geq 2$ !) » est écrit dans la marge gauche et relié par un
trait à « de façon unique »}

\emph{\struck{NB} Rescholie} \struck{Donc} Si $n \geq 2$, on a donc
\[
C_0 \xrightarrow{\ \approx\ } C'_0 \xrightarrow{\ \approx\ } C''_0
\]

\marginal{$n \geq 2$}
Formes quadratiques \struck{\ill{}} invariantes sur $V$ \struck{si \ill{}}
Celles qui sont $\neq 0$ sur tte fibre correspondent 1-1 aux hyperquadriques
\add{$Q_C$} de $C = P(V)$, invariantes par $\varphi^{C}_{\mathfrak{G}}$.
Supposons ici \emph{$\lambda_i = -1$} $\forall\, i$ (condition \struck{\ill{}}
\uncertain{nécessaire} pour qu'il existe une $f$ \struck{\ill{}} inv.\ non
\uncertain{dég.}\ $f_V(a_i)$ inv.\ \ldots) Il faut écrire
\marginal{\uncertain{dans} le \uncertain{groupe} $\mathfrak{G}$ \uncertain{ou}
$G$ \ldots\ les relations $\tau_i^{2} = 1$}
\[
f_V(a_i) a_i^{\vee} + \psi_V(a_i) = 0 \qquad 0 \leq i \leq n
\]
\note{sous $f_V(a_i)$, un signe $=$ vertical renvoie à $c_i$}
i.e.
\[
\varphi_V(a_i, a_j) = -c_i \langle a_j, a_i^{\vee} \rangle \qquad
0 \leq i, j \leq n
\]
donc
\note{la formule est précédée d'un repère illisible, peut-être « a) » ou un petit cercle}
\[
\varphi_V(a_i, a_j) = 0 \qquad i, j \
\text{non consécutifs.}
\]

\page{17}
\struck{b) $\varphi_V(a_i, a_i) = c_i \langle a_i, a_i^{\vee} \rangle$}
\note{ligne biffée ; sous ses deux membres, les accolades portent $-2c_i$ et
$-2$}

\[
\text{b)} \quad \underset{\substack{\Vert \\ 2 f(a_i) = 2c_i}}{\varphi_V(a_i, a_i)}
= -c_i \underbrace{\langle a_i, a_i^{\vee} \rangle}_{-2}
\qquad 0 \leq i \leq n \quad \text{OK}
\]
\[
\text{c)} \quad \varphi_V(a_i, a_{i+1}) = -c_i \underbrace{\langle a_{i+1}, a_i^{\vee} \rangle}_{\beta_i}
= -c_{i+1} \underbrace{\langle a_i, a_{i+1}^{\vee} \rangle}_{?}
\qquad 0 \leq i \leq n-1
\]
\note{l'accolade sous le dernier crochet porte une valeur surchargée,
illisible (rendue ici par un point d'interrogation) ; d'après la page 15, $\langle a_i, a_{i+1}^{\vee} \rangle = 1$}
i.e.
\[
\left\lbrace
\begin{array}{l}
\varphi_V(a_i, a_{i+1}) = -c_{i+1} \\
c_{i+1} = \beta_i c_i
\end{array}
\right.
\]
\[
c_i = (\beta_0 \cdots \beta_{i-1}) c_0 \qquad 1 \leq i \leq n
\]
\note{la borne supérieure de l'intervalle est d'une lecture incertaine}
donc $f_V$ multiple de
\[
\begin{aligned}
\underbrace{f_V^{0}}\Bigl(\sum_{0}^{n} x_i a_i\Bigr)
&= \sum_{i=0}^{n} (\beta_0 \cdots \beta_{i-1}) x_i^{2}
+ \sum_{0 \leq i \leq n-1} (\beta_0 \cdots \beta_i) x_i x_{i+1} \\
&= x_0^{2} + \beta_0 x_1^{2} + (\beta_0 \beta_1) x_2^{2} + \cdots
+ (\beta_0 \beta_1 \cdots \beta_{n-1}) x_n^{2} \\
&\quad + \beta_0 x_0 x_1 + \beta_0 \beta_1 x_1 x_2 + \cdots
+ (\beta_0 \cdots \beta_{n-1}) x_{n-1} x_n \\
&= x_0^{2} + \beta_0 x_1 (x_0 + x_1) + (\beta_0 \beta_1) x_2 (x_1 + x_2) + \cdots \\
&\quad + (\beta_0 \cdots \beta_{n-1}) x_n (x_{n-1} + x_n).
\end{aligned}
\]
\marginal{pour la suite noter $f_V$}
\note{les termes rectangles sont écrits avec le signe $+$. Or les relations
c) ci-dessus donnent $\varphi_V(a_i, a_{i+1}) = -c_{i+1} = -(\beta_0 \cdots
\beta_i) c_0$, ce qui, pour $c_0 = 1$, conduirait au coefficient $-(\beta_0
\cdots \beta_i)$ devant $x_i x_{i+1}$ ; ou bien un signe s'est perdu, ou bien
la convention (signe des $\beta_i$, normalisation de $\varphi_V$) n'est pas
celle qui est supposée ici. La page ne tranche pas. Le regroupement de la
dernière égalité est conforme à l'avant-dernière}

\emph{Donc} dans $C$ il y a \uncertain{une seule} hyperquadrique $Q_C$
invariante par $\mathfrak{G}$, et on a donc en fait
\[
G \longrightarrow \operatorname{Aut}(C, Q_C)
\]
\struck{$G \to SO(V, f_V)$}
\note{sous la dernière flèche, une seconde ligne, reliée à la première par
une petite flèche verticale, est biffée ; on y lit
$G \to SO(V, f_V)$, lecture incertaine}

\page{18}
Ceci nous amène à introduire

$\boxed{C'''_0}$ = cat.\ des \struck{triples} \add{quadruples}
$(C, Q_C, \varphi^{C}_{G}, h_0)$,
\marginal{NB $n \geq 2$}
$C$ fibré projectif de rg $n$, $Q_C$ hyperquadrique dans $C$,
$\varphi^{C}_{G} : G \to \operatorname{Aut}(C, Q_C)$, $h_0$ section de $C$,
satisfaisant les conditions
\begin{itemize}
\item[a)] les $\tau^{C}_{i} \overset{\mathrm{def}}{=} \varphi^{C}_{G}(\tau_i)$
sont des ps.\ réfl., strictes si $1 \leq i \leq n$, et associées à $h_0$ si
$i = 0$ (\uncertain{peut} être triviale si $i = 0$)
$\Longrightarrow$ $h_i$ $(1 \leq i \leq n)$, $H_i$ $(1 \leq i \leq n)$
\item[b)] Les $h_i$ lin.\ indép.\ en chaque fibre
\item[c)] $h_i \notin H_{i+1}$ en tte fibre $(0 \leq i \leq n-1)$
\end{itemize}

\marginal{NB $\tau^{C}_{0}$ est strict.\ ortho.\ sss c'est une réflexion
stricte, \uncertain{ce qui} \uncertain{p.} sss \struck{$\beta_0 \neq 0$}
i.e.\ sss $\forall\, s$, $\beta_0(s) \neq 0$ ou \uncertain{$k(s)$} de car.\
$\neq 2$}

\struck{et les $\tau_0$ stricts (\uncertain{\ill{}})} \struck{\emph{NB} si
$Q_C$ est lisse, alors \ill{} équivalent : les $H_i$ \ill{} \ill{}}
\note{deux lignes biffées d'un même trait sinueux}

\emph{NB} Supposons les $\tau_i$ stricts \uncertain{orthogonaux}
\add{dans $C$ \ill{}}, i.e.\ les $\beta_i$ inv.\ (\uncertain{il en est} ainsi
si $Q_C$ est lisse \ldots) Alors la condition c) équivaut à
\begin{itemize}
\item[c')] $(\tau_i \tau_{i+1})^{2} = \mathrm{id}$ dans chaque fibre, i.e.\
$\tau_i$ ne commute pas à $\tau_{i+1}$ en chaque fibre \quad
$0 \leq i \leq n-1$
\end{itemize}
\note{la page porte bien « $= \mathrm{id}$ » ; la paraphrase qui suit (« ne
commute pas ») et la condition c) demandent $(\tau_i \tau_{i+1})^{2} \neq
\mathrm{id}$, comme à la page 3 (« $\rho_i^{2} \neq \mathrm{id}$ en tte
fibre »). Probable lapsus, non corrigé}

La condition b) s'énonce ainsi
\begin{itemize}
\item[b')] Pour tt $s \in S$, la rep.\ \add{de $G$} induite sur $C_s$ est
irréductible [\struck{\ill{}} \uncertain{suffisance} claire, pour la
nécessité il faut travailler un peu \ldots]
\end{itemize}
\marginal{?}
\note{le point d'interrogation est dans la marge gauche, en regard de la
parenthèse carrée ; le reste de la page est blanc}

\page{19}
A) Cas $n$ pair $\geq 2$, donc $C$ de rang $n$ pair aussi, donc $V$ de rang
impair, donc
\[
O(f_V) \supset SO(f_V) \times \mu_2, \quad \text{et}
\]
\[
SO(f_V) \hookrightarrow \operatorname{Aut}(C, Q_C)
\]
En fait, l'hom.
\[
G \xrightarrow{\ \varphi^{C}_{G}\ } \operatorname{Aut}(C, Q_C)
\]
se factorise à travers $SO(f_V)$
\[
G \xrightarrow{\ \psi^{V}_{G}\ } SO(f_V)
\]
On notera que
\[
\psi^{V}_{G}(\tau_i) = -\varphi^{V}_{G}(\tau_i) = -\tau^{V}_{i}
\]
\note{le $\psi^{V}_{G}$ du premier membre est surchargé}

\emph{Ex} \emph{$n = 2$}, \struck{alors} Supposons $f_V$ \struck{lisse}
\uncertain{lisse} i.e.\ $Y = Q_C$ une fibrée en droites
\uncertain{projectives} i.e.\ $\beta_0, \beta_1$ et
$\beta_0 + \beta_1 - 4 = \alpha_0 + \alpha_1$ inv., alors
\[
\underset{\substack{\Vert \\ \alpha_0 + 2}}{\beta_0},\
\underset{\substack{\Vert \\ \alpha_1 + 2}}{\beta_1}
\]
\[
SO(f_V) \xrightarrow{\ \sim\ } \operatorname{Aut}(Q_C),
\]
donc la donnée de $\varphi^{C}_{G}$ ou $\psi^{V}_{G}$ équivaut à celle des
fibrés en droites projectives $Y$ $(= Q_C)$, et de l'opération de $G$ dessus
via $\tau^{Y}_{0}, \tau^{Y}_{1}, \tau^{Y}_{2}$. Traduisons les conditions
explicitées plus haut sur $\varphi^{C}_{G}$ en termes de $\varphi^{Y}_{G}$
(NB ici $\beta_0$ inv.\ donc $\tau_0$ réflexion stricte, donc $h_0$ s'en
déduit, \uncertain{est} \uncertain{un} $2$-\uncertain{rejet} \ill{} $C$ de
$Y$)
\note{avec $\beta_i = \alpha_i + 2$, on a bien $\beta_0 + \beta_1 - 4 =
\alpha_0 + \alpha_1$. La dernière ligne, serrée au bas de la feuille, est
d'une lecture très incertaine}

\page{20}
\begin{itemize}
\item[1)] $\forall\, 0 \leq i \leq 2$, $\tau^{C}_{i}$ est $\neq 1$ en chaque
fibre [$\Longrightarrow$ $h_0, h_1, h_2$ \quad $H_0, H_1, H_2$ dans $C$]
\item[2)] $h_0, h_1, h_2$ lin.\ indép.\ en chaque fibre
\struck{(\uncertain{à expliciter} \ill{} en termes de \ill{})}
$\Longleftrightarrow$ $\forall$ tte fibre \struck{\ill{}} $Y_{\bar{s}}$, il
n'existe pas de diviseurs $\Delta_s$ dans $Y_s$ de degré 2, invariant par
$\varphi^{Y}_{G}$, i.e.\ pas de restriction à une structure soit de droite
affine, soit de \struck{droite affine} torseur sous $G_m$ \ldots\ (ou
\uncertain{équivalent} sous $G_m$ \ldots)
\marginal{c'est sans doute ce qui exclut le cas « diédral »}
\item[3)] $\tau^{C}_{0} \tau^{C}_{1} \overset{\mathrm{def}}{=} \rho_f$ et
$\tau^{C}_{1} \tau^{C}_{2} \overset{\mathrm{def}}{=} \rho_s$ sont tels que
\[
\rho_f^{2} \neq \mathrm{id}_{Y}, \quad \rho_s^{2} \neq \mathrm{id}_{Y}
\quad \text{sur chaque fibre.}
\]
(cela exprime la non commutation de $\tau_0, \tau_1$ et de $\tau_1, \tau_2$,
sur chaque fibre
\end{itemize}
\note{au-dessus de « à expliciter », une surcharge biffée, illisible. Les
indices $f$ et $s$ de $\rho$ sont d'une lecture incertaine ; le premier
$\mathrm{id}$ porte un indice surchargé. La parenthèse finale n'est pas
refermée ; le reste de la page est blanc}

\end{document}
